\documentclass[12pt, titlepage]{article}

\usepackage{booktabs}
\usepackage{tabularx}
\usepackage{hyperref}
\hypersetup{
    colorlinks,
    citecolor=blue,
    filecolor=black,
    linkcolor=red,
    urlcolor=blue
}
\usepackage[round]{natbib}

\input{../Comments.text}
\renewcommand{\authornote}[3]{}
\renewcommand{\todo}[1]{}
\input{../Common.text}

\begin{document}

\title{System Verification and Validation Plan for \progname{}} 
\author{\authname}
\date{\today}
	
\maketitle

\pagenumbering{roman}

\section*{Revision History}

\begin{tabularx}{\textwidth}{p{3cm}p{2cm}X}
\toprule {\bf Date} & {\bf Version} & {\bf Notes}\\
\midrule
Oct 26 & 1.0 & Added system tests for authentication and task creation\\
Date 2 & 1.1 & Notes\\
\bottomrule
\end{tabularx}

~\\
\wss{The intention of the VnV plan is to increase confidence in the software.
However, this does not mean listing every verification and validation technique
that has ever been devised.  The VnV plan should also be a \textbf{feasible}
plan. Execution of the plan should be possible with the time and team available.
If the full plan cannot be completed during the time available, it can either be
modified to ``fake it'', or a better solution is to add a section describing
what work has been completed and what work is still planned for the future.}

\wss{The VnV plan is typically started after the requirements stage, but before
the design stage.  This means that the sections related to unit testing cannot
initially be completed.  The sections will be filled in after the design stage
is complete.  the final version of the VnV plan should have all sections filled
in.}

\newpage

\tableofcontents

\listoftables
\wss{Remove this section if it isn't needed}

\listoffigures
\wss{Remove this section if it isn't needed}

\newpage

\section{Symbols, Abbreviations, and Acronyms}

\renewcommand{\arraystretch}{1.2}
\begin{tabular}{l l} 
  \toprule		
  \textbf{symbol} & \textbf{description}\\
  \midrule 
  T & Test\\
  \bottomrule
\end{tabular}\\

\wss{symbols, abbreviations, or acronyms --- you can simply reference the SRS
  \citep{SRS} tables, if appropriate}

\wss{Remove this section if it isn't needed}

\newpage

\pagenumbering{arabic}

This document ... \wss{provide an introductory blurb and roadmap of the
  Verification and Validation plan}

\section{General Information}

\subsection{Summary}

\wss{Say what software is being tested.  Give its name and a brief overview of
  its general functions.}

\subsection{Objectives}

\wss{State what is intended to be accomplished.  The objective will be around
  the qualities that are most important for your project.  You might have
  something like: ``build confidence in the software correctness,''
  ``demonstrate adequate usability.'' etc.  You won't list all of the qualities,
  just those that are most important.}

\wss{You should also list the objectives that are out of scope.  You don't have 
the resources to do everything, so what will you be leaving out.  For instance, 
if you are not going to verify the quality of usability, state this.  It is also 
worthwhile to justify why the objectives are left out.}

\wss{The objectives are important because they highlight that you are aware of 
limitations in your resources for verification and validation.  You can't do everything, 
so what are you going to prioritize?  As an example, if your system depends on an 
external library, you can explicitly state that you will assume that external library 
has already been verified by its implementation team.}

\subsection{Custom Documents}

\wss{Summarize the CDs (Custom Documents) that were tackled by this project.  CDs 
can include usability testing, code walkthroughs, user documentation, formal
proof, GenderMag personas, Design Thinking, etc.  CDs should have already
been approved by the course instructor as included in your problem statement.
You can use a pull request to update your CDs (in TeamComposition.csv or
Repos.csv) if your plan changes as a result of the VnV planning exercise.}

\subsection{Relevant Documentation}

\wss{Reference relevant documentation.  This will definitely include your SRS
  and your other project documents (design documents, like MG, MIS, etc).  You
  can include these even before they are written, since by the time the project
  is done, they will be written.  You can create BibTeX entries for your
  documents and within those entries include a hyperlink to the documents.}

\citet{SRS}

\wss{Don't just list the other documents.  You should explain why they are relevant and 
how they relate to your VnV efforts.}

\section{Plan}

\wss{Introduce this section.  You can provide a roadmap of the sections to
  come.}

\subsection{Verification and Validation Team}

\wss{Your teammates.  Maybe your supervisor.
  You should do more than list names.  You should say what each person's role is
  for the project's verification.  A table is a good way to summarize this information.}

\subsection{SRS Verification}

\wss{List any approaches you intend to use for SRS verification.  This may
  include ad hoc feedback from reviewers, like your classmates (like your
  primary reviewer), or you may plan for something more rigorous/systematic.}

\wss{If you have a supervisor for the project, you shouldn't just say they will
read over the SRS.  You should explain your structured approach to the review.
Will you have a meeting?  What will you present?  What questions will you ask?
Will you give them instructions for a task-based inspection?  Will you use your
issue tracker?}

\wss{Maybe create an SRS checklist?}

\subsection{Design Verification}

\wss{Plans for design verification}

\wss{The review will include reviews by your classmates}

\wss{Create a checklists?}

\subsection{Verification and Validation Plan Verification}

\wss{The verification and validation plan is an artifact that should also be
verified.  Techniques for this include review and mutation testing.}

\wss{The review will include reviews by your classmates}

\wss{Create a checklists?}

\subsection{Implementation Verification}

\wss{You should at least point to the tests listed in this document and the unit
  testing plan.}

\wss{In this section you would also give any details of any plans for static
  verification of the implementation.  Potential techniques include code
  walkthroughs, code inspection, static analyzers, etc.}

\wss{The final class presentation in CAS 741 could be used as a code
walkthrough.  There is also a possibility of using the final presentation (in
CAS741) for a partial usability survey.}

\subsection{Automated Testing and Verification Tools}

\wss{What tools are you using for automated testing.  Likely a unit testing
  framework and maybe a profiling tool, like ValGrind.  Other possible tools
  include a static analyzer, make, continuous integration tools, test coverage
  tools, etc.  Explain your plans for summarizing code coverage metrics.
  Linters are another important class of tools.  For the programming language
  you select, you should look at the available linters.  There may also be tools
  that verify that coding standards have been respected, like flake9 for
  Python.}

\wss{If you have already done this in the development plan, you can point to
that document.}

\wss{The details of this section will likely evolve as you get closer to the
  implementation.}

\subsection{Software Validation}

\wss{If there is any external data that can be used for validation, you should
  point to it here.  If there are no plans for validation, you should state that
  here.}

\wss{You might want to use review sessions with the stakeholder to check that
the requirements document captures the right requirements.  Maybe task based
inspection?}

\wss{For those capstone teams with an external supervisor, the Rev 0 demo should 
be used as an opportunity to validate the requirements.  You should plan on 
demonstrating your project to your supervisor shortly after the scheduled Rev 0 demo.  
The feedback from your supervisor will be very useful for improving your project.}

\wss{For teams without an external supervisor, user testing can serve the same purpose 
as a Rev 0 demo for the supervisor.}

\wss{This section might reference back to the SRS verification section.}

\section{System Tests}

\wss{There should be text between all headings, even if it is just a roadmap of
the contents of the subsections.}

This section describes the tests that will be performed to verify that the entire integrated
system meets the functional and non-functional requirements specified in the \href{../SRS/SRS.pdf}{SRS}. The
system test cases below are be descrbed by the following characteristics:
\begin{itemize}
  \item Control: Is this test manual or automatic?
  \item Initial State: What is the state of the system before the test is performed?
  \item Input: What is the input to the system for this test?
  \item Output: What is the expected output of the system for this test?
  \item Test Case Derivation: What FRs or NFRs does this test case cover?
  \item How test will be performed: What is the abstract procedure for performing the test?
\end{itemize}

\subsection{Tests for Functional Requirements}

\wss{Subsets of the tests may be in related, so this section is divided into
  different areas.  If there are no identifiable subsets for the tests, this
  level of document structure can be removed.}

\wss{Include a blurb here to explain why the subsections below
  cover the requirements.  References to the SRS would be good here.}

\subsubsection{Accounts and Authentication}

These tests verify the account authentication requirements in the
\href{../SRS/SRS.pdf}{SRS}, including authentication with a compatible Web3
wallet extension and an embedded wallet.
		
\paragraph{Authentication Using a Web3 Wallet Extension}

\begin{enumerate}

\item{ST-AUTH-EXT-1\\}

Control: Manual

Initial State: The application is running in a supported browser on its public
page. A compatible Web3 wallet extension is installed and active. Test
profiles \texttt{EXT-USER-01} and \texttt{EXT-USER-02} are associated with
wallet fixtures \texttt{EXT-WALLET-01} and \texttt{EXT-WALLET-02}; neither
user is signed in.

Input:
\begin{itemize}
\item Connect \texttt{EXT-WALLET-01} and approve the verification-message
signature requested by the application.
\item Connect \texttt{EXT-WALLET-02} and approve its verification-message
signature.
\item Connect \texttt{EXT-WALLET-01} and reject the signature request.
\end{itemize}

Output:
\begin{itemize}
\item For the approved signature from \texttt{EXT-WALLET-01}, the user is
authenticated as \texttt{EXT-USER-01}, and that user's profile and
authenticated features are available.
\item For the approved signature from \texttt{EXT-WALLET-02}, the user is
authenticated as \texttt{EXT-USER-02}, and that user's profile and
authenticated features are available.
\item The user remains unauthenticated, no profile or authenticated features become available.
\end{itemize}

Test Case Derivation: FR1 requires the authentication
response to verify ownership of the connected wallet address.

How test will be performed: For each wallet, start from a signed-out session,
connect the specified extension account, and approve the verification-message
signature request. Confirm that \texttt{EXT-WALLET-01} opens
\texttt{EXT-USER-01}'s profile and \texttt{EXT-WALLET-02} opens
\texttt{EXT-USER-02}'s profile. Confirm that the user is granted access to the
authenticated features. Repeat the process and
reject the signature request; confirm that the user remains signed out and
cannot access authenticated features.

\item{ST-AUTH-EXT-2\\}

Control: Manual

Initial State: The application is running on its public page in a supported
browser. \texttt{EXT-WALLET-03} is available in a
compatible browser extension and is not yet associated with a DAIIM user
profile. The wallet owner is not signed in.

Input:
\begin{itemize}
\item Connect \texttt{EXT-WALLET-03} and approve the verification-message
signature.
\item Sign out of the resulting account.
\item Reconnect \texttt{EXT-WALLET-03} and approve the verification-message
signature again.
\end{itemize}

Output:
\begin{itemize}
\item After the first approved signature, the user is authenticated and a
profile is associated with \texttt{EXT-WALLET-03}. The system grants the user
access to authenticated features.
\item Sign-out ends the authenticated session; the user cannot access the
profile or authenticated features while signed out.
\item After the second approved signature, the same profile associated with
\texttt{EXT-WALLET-03} is displayed. No second profile is created for the
wallet. The system grants the user access to authenticated features.
\end{itemize}

Test Case Derivation: FR1 requires the system to
verify ownership of the wallet address and associate that verified address
with the same user profile on every successful authentication, including
after sign-out. The second successful authentication must restore
the profile created during the first sign-in.

How test will be performed: Record the number of DAIIM profiles associated
with \texttt{EXT-WALLET-03}. Connect the wallet and approve the
verification-message signature. Record the displayed profile identifier and
confirm that the wallet is associated with that profile. Sign out and verify
that the profile and authenticated features are no longer accessible. Sign in
again with the same wallet and approve the signature. Verify that the displayed
profile identifier matches the recorded identifier and that the number of
profiles associated with the wallet has not increased.

\end{enumerate}

\paragraph{Authentication Using an Embedded Wallet}

\begin{enumerate}

\item{ST-AUTH-EMB-1\\}

Control: Automatic

Initial State: The application is running in a test environment with the
embedded-wallet provider mocked. Identities \texttt{new.email@example.test},
\texttt{+15555550101}, and social-account subject \texttt{social-new-01} have
no associated DAIIM account or embedded wallet. Configure the provider mock
to return successful authentication and wallet creation for these identities,
with distinct addresses \texttt{EMB-WALLET-NEW-01},
\texttt{EMB-WALLET-NEW-02}, and \texttt{EMB-WALLET-NEW-03}, respectively.
The malformed email \texttt{not-an-email} should be rejected by the system
before the provider is called.

Input:
\begin{itemize}
\item Submit \texttt{new.email@example.test} using the email sign-in method.
\item Submit \texttt{+15555550101} using the phone sign-in method.
\item Submit social-account subject \texttt{social-new-01} using the supported
social sign-in method.
\item Submit malformed email \texttt{not-an-email} using the email sign-in
method.
\end{itemize}

Output:
\begin{itemize}
\item For each of the three supported identities, the system creates
one embedded wallet associated with the user's account, displays its
corresponding address and a creation confirmation, authenticates the user, and
makes the associated profile available.
\item For \texttt{not-an-email}, the system reports an input error, displays no
wallet-creation confirmation or address, and does not authenticate the user or
create a wallet.
\end{itemize}

Test Case Derivation: FR2 requires successful creation
for supported identity input after authentication, with the public address and
confirmation as outputs. The system must reject malformed identity input
before wallet creation. FR1 requires the successful
wallet authentication to identify the user.

How test will be performed: For each valid test identity, call the
corresponding embedded-wallet sign-in method with the specified email
address, phone number, or social-account identity. For each successful case,
compare the displayed address with the configured mock address, verify that
the corresponding profile is authenticated, and confirm that exactly one
wallet was created. For the malformed email, enter it using the email sign-in
method and verify that the system reports an input error, does not call the
provider, and does not create a wallet or authenticated session.

\item{ST-AUTH-EMB-2\\}

Control: Automatic

Initial State: The application is running in a test environment, and existing
accounts \texttt{EMB-RETURN-01}, \texttt{EMB-RETURN-02}, and
\texttt{EMB-RETURN-03} are signed out. Their embedded wallet addresses are
respectively \texttt{EMB-WALLET-RETURN-01},
\texttt{EMB-WALLET-RETURN-02}, and \texttt{EMB-WALLET-RETURN-03}. Configure
the embedded-wallet provider mock to return successful authentication
responses with those address mappings. Configure it to return an
authentication failure for \texttt{locked.user@example.test}, which has no
associated account.

Input:
\begin{itemize}
\item Enter \texttt{returning.email@example.test} using the email sign-in
method.
\item Enter \texttt{+15555550102} using the phone sign-in method.
\item Provide social-account subject \texttt{social-return-03} using the
supported social sign-in method.
\item Enter \texttt{locked.user@example.test} using the email sign-in method.
\end{itemize}

Output:
\begin{itemize}
\item For \texttt{returning.email@example.test}, the system authenticates
\texttt{EMB-RETURN-01} and displays its existing profile and
\texttt{EMB-WALLET-RETURN-01}.
\item For \texttt{+15555550102}, the system authenticates
\texttt{EMB-RETURN-02} and displays its existing profile and
\texttt{EMB-WALLET-RETURN-02}.
\item For social-account subject \texttt{social-return-03}, the system
authenticates \texttt{EMB-RETURN-03} and displays its existing profile and
\texttt{EMB-WALLET-RETURN-03}.
\item For \texttt{locked.user@example.test}, the system reports authentication
failure and does not create an authenticated session or grant access to a
profile.
\end{itemize}

Test Case Derivation: FR1 requires each verified
wallet address to remain associated with the same user profile after sign-out.
Each successful response must restore the mapped existing profile rather
than create another account. The failed authentication response does not verify
wallet ownership, so it must not grant access.

How test will be performed: Call the application's
embedded-wallet authentication entry point with the sign-in method and parameters from each returning-user credentials. Configure the provider
mock to return that fixture's existing wallet authentication response. Assert
that each call returns an authenticated session for the user profile. Then
call the same entry point with the sign-in method and identity parameters from
the authentication-failure credentials. Assert that the call returns no authenticated
session, grants no profile access and displays the appropriate error message.

\end{enumerate}

\subsubsection{Inference Task Creation}

This test verifies the end-to-end task creation flow, including validation of
the task definition, user authorization, and confirmation that the task-creation
program instruction succeeds on the blockchain.

\paragraph{Create and Validate Inference Tasks}

\begin{enumerate}

\item{ST-TASK-CREATE-1\\}

Control: Automatic

Initial State: The system is running on Solana Devnet. The requester
wallet address fixture is authenticated, and the test harness is configured with its
wallet signer and sufficient test SOL to pay transaction fees. The unique task
names \texttt{Regression-Test-Task},
\texttt{Regression-Test-Min-Providers},
\texttt{Regression-Test-Invalid-Schema}, and
\texttt{Rejected-Authorization-Task} are not already in use.

Input:
\begin{itemize}
\item A valid numerical regression task definition named
\texttt{Regression-Test-Task}, including a description, numerical input and
output schemas, an evaluation metric, a minimum provider count of 2 or more,
reputation tolerance, acceptance period, and maximum prediction time.
\item The valid definition named
\texttt{Regression-Test-Min-Providers}, with its minimum provider count
changed to 1.
\item The valid definition named \texttt{Regression-Test-Invalid-Schema},
with an input or output schema containing non-numerical values incompatible
with a numerical regression task.
\item The valid definition named \texttt{Rejected-Authorization-Task}.
\end{itemize}

Output:
\begin{itemize}
\item For the valid definition with authorization approved, the task is created
with the submitted definition, and the corresponding Solana Devnet transaction
succeeds with the task-creation instruction.
\item For the definition with fewer than 2 minimum providers, the system
rejects the definition, and no task is created.
\item For the definition with a schema incompatible with numerical regression,
the system rejects the definition, and no task is created.
\item For the valid definition with authorization rejected, no task is created
and no successful task-creation transaction occurs.
\end{itemize}

Test Case Derivation: FR6 requires a task name,
description, type, input schema, and output schema; the incompatible-schema
case checks that the definition matches the numerical regression task.
FR8 requires at least 2 providers, so a value of 1
must be rejected. FR5 permits authenticated users
to create tasks, while FR4 requires wallet
approval before an on-chain state change. Therefore, the valid task should be
created only after approval, and rejecting authorization should leave it
uncreated. FR50 and FR51
require a final transaction status and a transaction identifier that can be
used to verify the successful instruction.

How test will be performed: Submit the valid
definition to the task-creation flow and assert that the confirmation details
match the submitted definition. Configure the test wallet signer to approve
the authorization request. After confirmation, query Solana Devnet using the
returned transaction identifier and assert that the transaction succeeded and
contains the task-creation instruction. Query the application to confirm that
\texttt{Regression-Test-Task} is available with the submitted definition.
Submit the minimum-provider and invalid-schema definitions; assert that each
is rejected before authorization and no task is created. Finally, submit the
valid definition named \texttt{Rejected-Authorization-Task}, configure the
signer to reject authorization, and assert that no task-creation instruction
is submitted and no task is created.
\end{enumerate}

\subsubsection{Provider Registration and Staking}

These tests verify the provider registration and staking requirements in the
\href{../SRS/SRS.pdf}{SRS}. They cover registering as an inference provider
for a task, locking the required minimum stake in escrow, unregistering when
eligible, and viewing provider information through the dashboard.

\paragraph{Registering as an Inference Provider}

\begin{enumerate}

\item{ST-PROV-1\\}

Control: Manual

Initial State: The application is running in a supported browser. A provider
account is authenticated and connected to a compatible wallet. The provider
is not registered for inference task \texttt{TASK-01}. The task requires a
minimum stake of tokens, and the provider's wallet contains sufficient
tokens and SOL to complete the transaction.

Input:
\begin{itemize}
\item Navigate to \texttt{TASK-01} and select the option to register as an
inference provider.
\item Review the required minimum stake of tokens and the associated
registration details.
\item Confirm the registration and approve the wallet transaction.
\end{itemize}

Output:
\begin{itemize}
\item The provider is successfully registered for \texttt{TASK-01}.
\item Exactly the minimum amount of tokens are locked in escrow as the provider's stake.
\item The provider is added to the task's list of registered providers and
receives the base reputation score for the task.
\item The system displays the transaction status and, once confirmed, the
transaction identifier.
\end{itemize}

Test Case Derivation: FR\_InferenceTaskRegistration requires authenticated
users to be able to register as inference providers by staking the required
minimum amount of tokens. FR\_InferenceTaskRegistrationStaking requires the
minimum stake to be locked in escrow. The transaction authorization and
confirmation behaviour is also covered by the related transaction
requirements.

How test will be performed: Authenticate a provider account and navigate to
\texttt{TASK-01}. Select the registration option and verify that the system
displays the required minimum stake. Confirm the registration and approve
the transaction through the connected wallet. After the transaction is
confirmed, inspect the task's registered-provider list and the provider's
account information. Verify that the provider is registered, has received
the base reputation score, and that exactly the minimum amount of tokens have been transferred
into escrow.

\item{ST-PROV-2\\}

Control: Manual

Initial State: The application is running in a supported browser. A provider
account is authenticated and connected to a compatible wallet. The provider
is not registered for \texttt{TASK-02}. The task requires a minimum stake of
tokens, but the provider's wallet contains fewer than 100 tokens.

Input:
\begin{itemize}
\item Navigate to \texttt{TASK-02} and select the option to register as an
inference provider.
\item Attempt to proceed with registration despite having an insufficient
token balance.
\end{itemize}

Output:
\begin{itemize}
\item The system prevents the provider from completing registration.
\item An error message indicates that the provider's token balance is
insufficient and identifies the required shortfall.
\item No stake is transferred into escrow, and the provider is not added to
the task's list of registered providers.
\end{itemize}

Test Case Derivation: FR\_InferenceTaskRegistration and
FR\_InferenceTaskRegistrationStaking require the provider to stake the
required minimum amount of tokens to register. The related balance-check
assumptions require the system to verify that sufficient funds are available
before processing the transaction.

How test will be performed: Authenticate a provider account with fewer than
100 tokens and navigate to \texttt{TASK-02}. Attempt to register for the task.
Verify that the system displays an insufficient-balance error, that no
transaction transfers tokens into escrow, and that the provider does not
appear in the task's registered-provider list.

\end{enumerate}

\paragraph{Unregistering from an Inference Task}

\begin{enumerate}

\item{ST-PROV-3\\}

Control: Manual

Initial State: The application is running in a supported browser. A provider
account is authenticated and registered for \texttt{TASK-03}. The provider
has no active accepted inference requests for this task. The provider's
current reputation score is at or above the task's reduced-stake-return
threshold, and the provider has a remaining stake of tokens.

Input:
\begin{itemize}
\item Navigate to the provider's registration details for
\texttt{TASK-03}.
\item Select the option to unregister from the task.
\item Review and confirm the unregistration.
\item Approve the wallet transaction, if prompted.
\end{itemize}

Output:
\begin{itemize}
\item The provider is successfully unregistered from \texttt{TASK-03}.
\item The provider is removed from the task's list of registered providers.
\item The provider's entire remaining stake of tokens is returned to
their token account.
\item The system displays the transaction status and, once confirmed, the
transaction identifier.
\end{itemize}

Test Case Derivation: FR\_ProviderWillingUnregistration requires the system
to allow providers to unregister when they have no active accepted inference
requests. FR\_ProviderUnregistrationStakeReturn requires the entire
remaining stake to be returned when the provider's reputation is at or above
the task's reduced-stake-return threshold.

How test will be performed: Authenticate the registered provider and verify
that they have no active accepted inference requests for
\texttt{TASK-03}. Select the unregistration option and confirm the action.
After the transaction is confirmed, inspect the provider's registration
status, the task's registered-provider list, and the provider's token
balance. Verify that the provider has been removed from the list and that
the entire remaining stake has been returned.

\item{ST-PROV-4\\}

Control: Manual

Initial State: The application is running in a supported browser. A provider
account is authenticated and registered for \texttt{TASK-04}. The provider
has accepted an inference request that is still active for this task.

Input:
\begin{itemize}
\item Navigate to the provider's registration details for
\texttt{TASK-04}.
\item Select the option to unregister from the task.
\item Attempt to confirm the unregistration.
\end{itemize}

Output:
\begin{itemize}
\item The system rejects the unregistration attempt because the provider has
an active accepted inference request.
\item An error message explains why the provider cannot unregister.
\item The provider remains registered for \texttt{TASK-04}, and their
remaining stake is unchanged.
\end{itemize}

Test Case Derivation: FR\_ProviderWillingUnregistration requires that a
provider must not have accepted any active inference requests for the task
before they can unregister.

How test will be performed: Authenticate a provider account that is
registered for \texttt{TASK-04} and has an active accepted inference request.
Attempt to unregister from the task. Verify that the system rejects the
attempt, displays an appropriate error message, maintains the provider's
registration, and does not modify the provider's remaining stake.

\end{enumerate}

\paragraph{Viewing the Provider Dashboard}

\begin{enumerate}

\item{ST-PROV-5\\}

Control: Manual

Initial State: The application is running in a supported browser. A provider
account is authenticated and registered for two inference tasks,
\texttt{TASK-05} and \texttt{TASK-06}. The provider has an existing reputation
score, staked token balance, and earned rewards recorded for each task.

Input:
\begin{itemize}
\item Navigate to the provider dashboard.
\item View the provider's information for \texttt{TASK-05} and
\texttt{TASK-06}.
\end{itemize}

Output:
\begin{itemize}
\item The dashboard displays the provider's reputation score, staked tokens,
and earned rewards for \texttt{TASK-05}.
\item The dashboard displays the provider's reputation score, staked tokens,
and earned rewards for \texttt{TASK-06}.
\item The displayed values correspond to the provider's recorded information
for each task.
\end{itemize}

Test Case Derivation: FR\_ProviderDashboard requires inference providers to
be able to view their reputation score, staked tokens, and rewards earned
for each inference task they are registered for.

How test will be performed: Authenticate the provider account and navigate
to the provider dashboard. Retrieve the provider's recorded reputation,
stake, and reward values for both tasks from the system's records. Compare
these values with the information displayed on the dashboard and verify
that the correct information is displayed for each task.

\end{enumerate}

\subsubsection{Inference Requests}

These tests verify the inference request requirements in the
\href{../SRS/SRS.pdf}{SRS}. They cover submitting requests for existing
inference tasks, validating request input data against the task's input
schema, authorizing payments into escrow, and broadcasting submitted
requests to registered providers.

\paragraph{Submitting an Inference Request}

\begin{enumerate}

\item{ST-REQ-1\\}

Control: Manual

Initial State: The application is running in a supported browser. A requester
account is authenticated and connected to a compatible wallet. An existing
inference task, \texttt{TASK-01}, is available and has at least two registered
providers. The task's input schema requires a numerical input field named
\texttt{input\_value}. The requester's wallet contains sufficient tokens and
SOL to submit the request.

Input:
\begin{itemize}
\item Navigate to \texttt{TASK-01} and select the option to submit an
inference request.
\item Provide valid input data containing \texttt{input\_value} with a
numerical value
\item Review the request details and confirm the submission.
\item Authorize the payment through the connected wallet.
\end{itemize}

Output:
\begin{itemize}
\item The system accepts the input data and successfully submits the
inference request for \texttt{TASK-01}.
\item The request is assigned a unique identifier and its submission status
is displayed to the requester.
\item The input data is stored securely for processing and is not recorded
directly on-chain.
\item The requester is informed that the request has been submitted and is
awaiting provider participation.
\end{itemize}

Test Case Derivation: FR\_RequestSubmission requires authenticated users to
be able to submit inference requests for existing tasks, with input data
conforming to the task's input schema. The related input-data requirements
specify that input data must be stored securely and not placed directly
on-chain.

How test will be performed: Authenticate a requester account and navigate
to \texttt{TASK-01}. Enter the valid input data and submit the request.
Authorize the transaction through the connected wallet. Verify that the
system creates a request with a unique identifier, displays the appropriate
submission status, and securely stores the input data. Inspect the on-chain
request record to confirm that the input data itself has not been recorded
on-chain.

\item{ST-REQ-2\\}

Control: Manual

Initial State: The application is running in a supported browser. A requester
account is authenticated and connected to a compatible wallet. An existing
inference task, \texttt{TASK-02}, is available and has at least two registered
providers. The task's input schema requires a numerical field named
\texttt{input\_value}.

Input:
\begin{itemize}
\item Navigate to \texttt{TASK-02} and select the option to submit an
inference request.
\item Provide input data with the required \texttt{input\_value} field
missing.
\item Repeat the submission attempt with \texttt{input\_value} provided as
a non-numerical value.
\end{itemize}

Output:
\begin{itemize}
\item The system rejects both invalid input submissions.
\item The system displays an error message identifying the missing or
invalid input field.
\item No inference request is submitted, no request identifier is created,
and no payment is transferred into escrow.
\end{itemize}

Test Case Derivation: FR\_RequestValidation requires the system to reject
inference requests whose input data does not conform to the task's input
schema. FR\_RequestSubmission requires submitted requests to include input
data that conforms to that schema.

How test will be performed: Authenticate a requester account and navigate
to \texttt{TASK-02}. Attempt to submit a request with the required input
field missing. Verify that the system rejects the request and identifies
the missing field. Repeat the test using a non-numerical value for
\texttt{input\_value}. Verify that the request is rejected again and that
neither attempt creates a request or transfers funds into escrow.

\end{enumerate}

\paragraph{Request Payment and Authorization}

\begin{enumerate}

\item{ST-REQ-3\\}

Control: Manual

Initial State: The application is running in a supported browser. A requester
account is authenticated and connected to a compatible wallet. An existing
inference task, \texttt{TASK-03}, is available and has at least two registered
providers. The task has a minimum payment requirement. The
requester's wallet contains sufficient tokens and SOL to authorize a
payment 

Input:
\begin{itemize}
\item Navigate to \texttt{TASK-03} and provide valid input data.
\item Select a payment amount of tokens.
\item Review the transaction details and confirm the request submission.
\item Approve the transaction through the connected wallet.
\end{itemize}

Output:
\begin{itemize}
\item The system accepts the payment amount because it exceeds the task's
minimum payment requirement.
\item The wallet prompts the requester to authorize a transaction
transferring the tokens into escrow.
\item After authorization and transaction confirmation, exactly those amount of tokens
are transferred into escrow.
\item The system displays the transaction status and, once confirmed, the
transaction identifier.
\item The amount transferred into escrow matches the amount authorized by
the requester.
\end{itemize}

Test Case Derivation: FR\_RequestPayment requires the requester to choose
a token payment greater than the minimum and authorize a wallet transaction
transferring that amount into escrow. The related transaction authorization
and preview requirements ensure that the requester can review and authorize
the transaction. SA\_ChargeEqualsAuthorized requires the amount charged to
match the amount authorized.

How test will be performed: Authenticate a requester account and navigate
to \texttt{TASK-03}. Enter valid request input and select a payment amount
for a set amount of tokens. Review the transaction preview and approve the transaction
through the connected wallet. After confirmation, inspect the requester's
token balance and the escrow account's balance. Verify that exactly the set amount of
tokens have been transferred into escrow and that the amount matches the
requester's authorization.

\end{enumerate}

\paragraph{Broadcasting Submitted Inference Requests}

\begin{enumerate}

\item{ST-REQ-4\\}

Control: Automatic

Initial State: The system is running in a test environment. An existing
inference task, \texttt{TASK-04}, has three registered providers:
\texttt{PROV-01}, \texttt{PROV-02}, and \texttt{PROV-03}. All three providers
are available to receive inference requests. The task has a configured
acceptance period and maximum prediction time.

Input:
\begin{itemize}
\item Submit a valid inference request for \texttt{TASK-04}.
\item Set the proposed payment amount to a number of tokens.
\item Configure the acceptance-period expiry time to 10 minutes after
submission.
\item Configure the maximum prediction time to 30 minutes.
\end{itemize}

Output:
\begin{itemize}
\item The system successfully submits the inference request.
\item The system broadcasts the request to all three registered providers.
\item Each provider receives the task definition, proposed payment amount,
acceptance-period expiry time, and maximum prediction time.
\item The broadcast information is consistent across all three providers
and matches the submitted request.
\end{itemize}

Test Case Derivation: FR\_RequestBroadcast requires the system to broadcast
the task definition, proposed payment amount, acceptance-period expiry time,
and maximum prediction time to all providers registered for the task after
an inference request is submitted.

How test will be performed: Configure \texttt{TASK-04} with the three
specified registered providers and submit a valid inference request with
the specified payment amount and timing values. Monitor the messages or
notifications received by each provider. Verify that all three providers
receive the broadcast and that the task definition, proposed payment,
acceptance-period expiry time, and maximum prediction time match the
submitted request.

\end{enumerate}


\subsection{Tests for Nonfunctional Requirements}

\wss{The nonfunctional requirements for accuracy will likely just reference the
  appropriate functional tests from above.  The test cases should mention
  reporting the relative error for these tests.  Not all projects will
  necessarily have nonfunctional requirements related to accuracy.}

\wss{For some nonfunctional tests, you won't be setting a target threshold for
passing the test, but rather describing the experiment you will do to measure
the quality for different inputs.  For instance, you could measure speed versus
the problem size.  The output of the test isn't pass/fail, but rather a summary
table or graph.}

\wss{Tests related to usability could include conducting a usability test and
  survey.  The survey will be in the Appendix.}

\wss{Static tests, review, inspections, and walkthroughs, will not follow the
format for the tests given below.}

\wss{If you introduce static tests in your plan, you need to provide details.
How will they be done?  In cases like code (or document) walkthroughs, who will
be involved? Be specific.}

\subsubsection{Area of Testing1}
		
\paragraph{Title for Test}

\begin{enumerate}

\item{test-id1\\}

Type: Functional, Dynamic, Manual, Static etc.
					
Initial State: 
					
Input/Condition: 
					
Output/Result: 
					
How test will be performed: 
					
\item{test-id2\\}

Type: Functional, Dynamic, Manual, Static etc.
					
Initial State: 
					
Input: 
					
Output: 
					
How test will be performed: 

\end{enumerate}

\subsubsection{Area of Testing2}

...

\subsection{Traceability Between Test Cases and Requirements}

\wss{Provide a table that shows which test cases are supporting which
  requirements.}

\section{Unit Test Description}

\wss{This section should not be filled in until after the MIS (detailed design
  document) has been completed.}

\wss{Reference your MIS (detailed design document) and explain your overall
philosophy for test case selection.}  

\wss{To save space and time, it may be an option to provide less detail in this section.  
For the unit tests you can potentially layout your testing strategy here.  That is, you 
can explain how tests will be selected for each module.  For instance, your test building 
approach could be test cases for each access program, including one test for normal behaviour 
and as many tests as needed for edge cases.  Rather than create the details of the input 
and output here, you could point to the unit testing code.  For this to work, you code 
needs to be well-documented, with meaningful names for all of the tests.}

\subsection{Unit Testing Scope}

\wss{What modules are outside of the scope.  If there are modules that are
  developed by someone else, then you would say here if you aren't planning on
  verifying them.  There may also be modules that are part of your software, but
  have a lower priority for verification than others.  If this is the case,
  explain your rationale for the ranking of module importance.}

\subsection{Tests for Functional Requirements}

\wss{Most of the verification will be through automated unit testing.  If
  appropriate specific modules can be verified by a non-testing based
  technique.  That can also be documented in this section.}

\subsubsection{Module 1}

\wss{Include a blurb here to explain why the subsections below cover the module.
  References to the MIS would be good.  You will want tests from a black box
  perspective and from a white box perspective.  Explain to the reader how the
  tests were selected.}

\begin{enumerate}

\item{test-id1\\}

Type: \wss{Functional, Dynamic, Manual, Automatic, Static etc. Most will
  be automatic}
					
Initial State: 
					
Input: 
					
Output: \wss{The expected result for the given inputs}

Test Case Derivation: \wss{Justify the expected value given in the Output field}

How test will be performed: 
					
\item{test-id2\\}

Type: \wss{Functional, Dynamic, Manual, Automatic, Static etc. Most will
  be automatic}
					
Initial State: 
					
Input: 
					
Output: \wss{The expected result for the given inputs}

Test Case Derivation: \wss{Justify the expected value given in the Output field}

How test will be performed: 

\item{...\\}
    
\end{enumerate}

\subsubsection{Module 2}

...

\subsection{Tests for Nonfunctional Requirements}

\wss{If there is a module that needs to be independently assessed for
  performance, those test cases can go here.  In some projects, planning for
  nonfunctional tests of units will not be that relevant.}

\wss{These tests may involve collecting performance data from previously
  mentioned functional tests.}

\subsubsection{Module ?}
		
\begin{enumerate}

\item{test-id1\\}

Type: \wss{Functional, Dynamic, Manual, Automatic, Static etc. Most will
  be automatic}
					
Initial State: 
					
Input/Condition: 
					
Output/Result: 
					
How test will be performed: 
					
\item{test-id2\\}

Type: Functional, Dynamic, Manual, Static etc.
					
Initial State: 
					
Input: 
					
Output: 
					
How test will be performed: 

\end{enumerate}

\subsubsection{Module ?}

...

\subsection{Traceability Between Test Cases and Modules}

\wss{Provide evidence that all modules have been considered.}
				
\bibliographystyle{plainnat}

\bibliography{../../refs/References}

\newpage

\section{Appendix}

This is where you can place additional information.

\subsection{Symbolic Parameters}

The definition of the test cases will call for SYMBOLIC\_CONSTANTS.
Their values are defined in this section for easy maintenance.

\subsection{Usability Survey Questions?}

\wss{This is a section that would be appropriate for some projects.}

\newpage{}

\section{Appendix --- Generative AI Use Disclosure}

\subsection{AI Tools Used}

An AI assistant using the Copilot SDK in Visual Studio Code was used to assist
with revisions to this document. The specific model used is GPT-6 Luna.

\subsection{How and Where Used}

The AI assistant was used to discuss, draft, and revise prose and requirements
for the verification and validation work. In the VnV Plan, it helped structure
system tests by functional area and draft and refine test cases for wallet
authentication, embedded-wallet creation, returning-user authentication, and
inference-task creation. Revisions focused on making inputs, expected outputs,
test procedures, and requirement derivations clearer and more specific.

The assistant also helped revise related SRS content: clarifying support for
embedded and browser-extension wallets, adding requirements for embedded-wallet
creation and persistent wallet-to-profile association, and adding a
conservative User Profile entry to the major data entities and data
dictionary. The team directed the scope and reviewed or adjusted the proposed
content.

\subsection{Verification of AI-Generated Content}

AI-assisted content in this document was checked against the SRS to confirm
that each test case addresses stated requirements and does not introduce
unapproved system behaviour or design decisions. Requirement references and
the expected outcomes derived from them were checked against the corresponding
SRS requirements and use cases. Test cases were reviewed for clear inputs,
expected results, and procedures, and the LaTeX structure and cross-references
were inspected after editing. Test fixtures, schema examples, and blockchain
verification details are proposed examples and must be reconciled with the
implemented system before execution. The team remains responsible for the
accuracy, completeness, and consistency of this document, regardless of
whether material was first drafted by a team member or with the help of an AI
tool.

\newpage{}

\section{Appendix --- Reflection}

\wss{This section is not required for CAS 741}

The information in this section will be used to evaluate the team members on the
graduate attribute of Lifelong Learning.

\input{../Reflection.text}

\begin{enumerate}
  \item What went well while writing this deliverable? 
  \item What pain points did you experience during this deliverable, and how
    did you resolve them?
  \item What knowledge and skills will the team collectively need to acquire to
  successfully complete the verification and validation of your project?
  Examples of possible knowledge and skills include dynamic testing knowledge,
  static testing knowledge, specific tool usage, Valgrind etc.  You should look to
  identify at least one item for each team member.
  \item For each of the knowledge areas and skills identified in the previous
  question, what are at least two approaches to acquiring the knowledge or
  mastering the skill?  Of the identified approaches, which will each team
  member pursue, and why did they make this choice?
\end{enumerate}

\end{document}